\subsection{Spinal Load Provider Split}
\label{sec:spn-12}

This section documents the separation of heuristic and model-based spinal
loads (issue \#12765, epic \#12735). Code:
\texttt{src/shared/python/injury/spinal\_load\_providers.py} and the provider
seam in \texttt{src/shared/python/injury/spinal\_load\_analysis.py}. The
\texttt{spine} package does not import \texttt{injury}.

\paragraph{Interface.}
\texttt{SpinalLoadProvider} exposes \texttt{provider\_id}, \texttt{is\_model\_based}
and \texttt{loads(time\_s, joint\_angles, joint\_velocities, joint\_torques,
levels)}, returning a \texttt{SpinalLoadSeries} in the ADR-0055 convention
(wrench on the superior vertebra by the inferior one, inferior ISB frame,
compression $+F_y$, AP shear $+F_x$, ML shear $+F_z$; N and N\,m).

\paragraph{Heuristic Provider.}
\texttt{HeuristicSpinalLoadProvider} (\texttt{provider\_id = heuristic\_v1},
\texttt{is\_model\_based = False}) wraps the unchanged analyzer closed forms.
For level $i$ with position ratio $p_i$, body mass $m$ and head-arms-trunk
ratio $0.678$, the mass above is $m_a = 0.678\,m\,(1-p_i)$ and
\begin{align}
  F_y &= m_a g\cos\phi + |\tau_\phi|/0.05, &
  F_x &= m_a g\sin\phi, \\
  F_z &= m_a g\sin\lambda + \tfrac12 m_a \omega^2 L p_i, &
  M_y &= |\tau_\rho|,
\end{align}
with lumbar flexion $\phi$, lateral bend $\lambda$, rotation rate $\omega$,
trunk length $L$ and the fixed 5\,cm moment arm. These are screening
estimates, not joint reactions. Every heuristic series carries
\texttt{solver = heuristic\_v1}, a \texttt{distribution\_source} stating
``heuristic'', \texttt{provider\_label = heuristic} and \texttt{model\_id =
n/a (heuristic)}, so it cannot be mistaken for a model load.

\paragraph{Model Provider.}
\texttt{ModelSpinalLoadProvider} wraps a precomputed series (the OpenSim
pipeline output, delivered only through the \texttt{SpinalLoadSeries} contract)
and returns it unchanged with provenance intact. A requested time base must
equal the series time base to $10^{-9}$\,s and requested levels must be present;
otherwise \texttt{ValueError}. Its \texttt{provider\_id} may not equal
\texttt{heuristic\_v1}.

\paragraph{Analyzer Behaviour.}
\texttt{SpinalLoadAnalyzer(provider=None)} defaults to the heuristic provider
and its outputs are unchanged (golden-tested to $10^{-12}$). Results expose
\texttt{provider\_id}, \texttt{is\_model\_based} and \texttt{risk\_assessed}.
With a model provider the segments come from the series (torsion is
$|M_y|$) and the heuristic risk bands are \emph{not} applied
(\texttt{risk\_assessed = False}, and every risk field reads \texttt{NOT\_ASSESSED}, never \texttt{SAFE}); the literature thresholds were derived for
the heuristic quantities. The spinal compression/shear terms in
\texttt{optimization/\_swing\_objectives.py} are labelled heuristic
(\texttt{SPINAL\_OBJECTIVE\_PROVIDER\_ID}); numerics are unchanged.

\paragraph{Limitations.}
The heuristic series maps only analyzed segment names that are valid spinal
levels; custom names raise from \texttt{loads()} but not from the default
\texttt{analyze()} path. No OpenSim computation occurs here.

\paragraph{Reproduction.}
\begin{verbatim}
python -m pytest tests/unit/injury tests/analysis/test_spinal_load_analysis.py \
  -q -n 0 -p no:cacheprovider -o addopts=""
\end{verbatim}
